% Section 19.1, from lectures/L19/appendix/a_system_verification.md.

\section{Verifying the whole system}
\label{c19:sec:verify}

Before any part of the transport can be trusted, the controller has to be shown working as a
\emph{system}. That does not mean one module against a stub, but two independent nodes on one wire
fighting over it the way real nodes do.

\subsection*{Modelling the bus}
Nothing in the design models the shared bus, and that is deliberate. \code{can_controller} drives a
pair \code{(tx_bus, bus_en)} and reads \code{rx_bus}, which is exactly what a real CAN transceiver
presents on its digital side (\chapref{c10:ch}). The wired-AND lives in the testbench, and on hardware
it lives in the transceiver.

\file{can_controller_tb.vhd} instantiates two nodes on one wire and combines them directly. Every
node's effective output is its driven value when \code{bus_en = '1'}, otherwise released
(\code{'1'}), and the bus is their AND, so it reads dominant the moment \emph{any} of the nodes pulls
low:

\begin{vhdlcode}
bus_line <= (a_tx_bus or not a_bus_en) and (b_tx_bus or not b_bus_en);
\end{vhdlcode}

That single line is the whole of \secref{c10:sec:can}'s ``any dominant beats all recessives'', and it
is what makes arbitration testable in simulation.

Worth working through rather than reading. Every node contributes \code{tx_bus or not bus_en}, which
is \code{1} as soon as it has released the bus, and the wire is the two contributions ANDed:

\begin{narrowtable}{llccccl}
  \thead{Node A} & \thead{Node B} & \thead{\code{a_tx_bus}} & \thead{\code{a_bus_en}} &
    \thead{\code{b_tx_bus}} & \thead{\code{b_bus_en}} & \thead{\code{bus_line}} \\
  \midrule
  released & released & \code{-} & \code{0} & \code{-} & \code{0} & \code{1} recessive \\
  driving dominant & released & \code{0} & \code{1} & \code{-} & \code{0} & \code{0} dominant \\
  released & driving dominant & \code{-} & \code{0} & \code{0} & \code{1} & \code{0} dominant \\
  driving dominant & driving dominant & \code{0} & \code{1} & \code{0} & \code{1} & \code{0}
    dominant \\
\end{narrowtable}

A released node's \code{tx_bus} is \code{-} because it genuinely does not matter: \code{not bus_en}
alone makes its contribution \code{1}, so a node that has let go cannot affect the wire whatever its
output happens to sit at. Across all sixteen input combinations the wire reads dominant exactly when
at least one node has \code{bus_en = '1'} with \code{tx_bus = '0'}, which is the rule expressed as
hardware.

Two rows carry the arbitration argument. \textbf{Row two is the contest itself:} A drives dominant, B
sends recessive and has therefore released, and the wire reads dominant. B, monitoring while it
transmits, sees that it sent recessive and that the bus reads dominant, and that is exactly the loss
condition \chapref{c17:ch} checks. Meanwhile A reads back precisely the bit it sent and notices
nothing.

\textbf{Row one is why the loser can leave.} Once B has stopped driving it contributes \code{1}
forever, so the wire is A's alone. The loser's retreat is invisible to the winner, and that is what
makes a collision a non-event rather than something both nodes have to recover from.

Row four is worth a look too: when both nodes drive dominant the wire is dominant and each reads back
what it sent, so neither can tell the other exists. That is the normal state for every dominant bit of
the identifier before the two IDs differ.

\subsection*{The two scenarios}
\begin{enumerate}
  \item \textbf{Reception.} Node A sends a 5-byte frame (ID \code{0x123}, data
    \code{11 22 F8 44 55}) while node B receives it. \code{0xF8} forces a stuff bit inside the data
    field, so both sides' CRC engines have to skip it. The testbench checks A's \code{tx_done} and
    B's \code{rx_valid}, that B reconstructed the ID, the DLC and \textbf{every} data byte at its
    fixed position, and that neither node raised \code{error}.
  \item \textbf{Arbitration.} Nodes A (ID \code{0x000}) and B (ID \code{0x001}) both request during
    the same cycle, differing only in the identifier's last bit. A has to win and complete normally,
    unaware that any contest took place; B has to detect the loss at that bit, raise \code{error} and
    stop driving. The check runs after A's frame has completed, so it also confirms that B's
    \code{error} is still readable then, rather than having been wiped by a re-entry into
    \code{STATE_START}.
\end{enumerate}

A transmitting node never receives its own frame (\code{rx_shift_reg} runs only while
\code{role = '0'}), so both scenarios genuinely need two nodes; neither could be tested with one node
alone.

The two one-cycle pulses the first scenario waits for, A's \code{tx_done} and B's \code{rx_valid}, are
\textbf{latched} in the testbench, since two independent state machines need not pulse at the same
edge, and a plain \code{wait until (a = '1' or b = '1')} would resume on the first and miss the
second.

\subsection*{What the testbench does not cover}
\label{c19:sub:gaps-tb}

Two passing scenarios are not the same thing as a proved design, and knowing where a testbench ends is
part of reading it. Five gaps are worth naming, since each is a place where this design is either
deliberately simplified or simply untested:
\begin{itemize}
  \item \textbf{Both nodes share one \code{clock} and one \code{reset_n}.} Real nodes run on
    independent oscillators and drift apart. Be precise about what that leaves untested, because it is
    not \code{resync} itself: both nodes pulse \code{resync} at every frame's start (\chapref{c16:ch}),
    and that pulse is what phases B's sample points against the arbitrary moment A began transmitting,
    so case 1 fails without it. What a shared clock can never exercise is \textbf{drift}, two timers
    starting in phase and wandering apart mid-frame, the case
    Exercise~\ref{c12:ex:drift} quantifies and the reason real CAN resynchronises at every
    recessive-to-dominant edge instead of once per frame. This is the largest gap, and closing it in
    simulation takes a second clock generator with a deliberate frequency deviation.
    \Secref{c19:sec:bringup} gets a frame across \textbf{two boards} with two separate crystals,
    finally genuinely independent clocks, with an honest caveat noted there.
  \item \textbf{No frame is ever corrupted.} Nothing injects a stuffing error, flips a bit on the wire
    or damages a CRC, so the receive side's \code{error} paths are exercised only at module level by
    \code{rx_shift_reg_tb} and \code{crc15_tb}, never end to end.
  \item \textbf{The ACK slot is driven but never read back.} A receiver with a valid CRC pulls it
    dominant (\chapref{c17:ch}), but the transmitter never checks it, so a frame nobody acknowledged
    completes exactly like one that was. That is a deliberate simplification
    (\secref{c19:sec:gaps}), and the testbench inherits it instead of catching it.
  \item \textbf{Arbitration is always two-way, and always lost at the same bit.} \code{0x000} against
    \code{0x001} differ in the identifier's last bit, so the loss is detected in the same state every
    run. Exercise~\ref{c19:ex:thirdnode} adds a third node and moves the losing bit earlier.
  \item \textbf{No node ever transmits again after an abort.} Node B loses arbitration in case 2 and
    is never asked to send another frame, so nothing here checks that a node recovers from an abort at
    all. That matters more than it sounds: \code{can_controller} clears \code{crc15} in
    \code{STATE_START} (\chapref{c16:ch}) precisely so that an abandoned frame cannot leave its residue
    behind, and this testbench would pass just as happily with that clear removed.
    Exercise~\ref{c19:ex:afterabort} is where you find out what it did.
\end{itemize}
